\documentclass[11pt,a4paper]{article}
\usepackage[T2A]{fontenc}
\usepackage[utf8]{inputenc}
\usepackage[russian,english]{babel}
\usepackage{amsmath,amssymb,amsthm,mathtools}
\usepackage{setspace}
\usepackage{enumitem}
\usepackage[hidelinks]{hyperref}
\setstretch{1.3}
\frenchspacing
\hyphenpenalty=10000

\newcommand{\head}[1]{\par\noindent\textbf{\underline{#1}}\quad}
\newcommand{\sheettitle}[4]{% title, subtitle, homework-flag, date
  \begin{center}
  {\huge\bfseries #1\par}\vspace{\dimexpr1.2em-1pt\relax minus 3.6pt}
  {\Large #2\par}\vspace{\dimexpr1em+1pt\relax minus 1.2pt}
  \if\relax\detokenize{#3}\relax\else{\Large #3\par}\vspace{\dimexpr1.15em-1pt\relax minus 3.6pt}\fi
  {\foreignlanguage{russian}{#4}\par}
  \end{center}\vspace{0.5pt minus 1.2pt}}
\newcommand{\src}[1]{\unskip\nobreak\hfil\penalty50\hskip1em\hbox{}\nobreak\hfill(#1){\parfillskip=0pt\par}}
\DeclareMathOperator{\tr}{tr}
\DeclareMathOperator{\adj}{adj}
\newcommand{\T}{\mathsf T}
\DeclareMathOperator{\rank}{rank}
\begin{document}
\sheettitle{Linear algebra -- 9}{Matrices over $\mathbb Z$ and over a finite field}{Homework}{12 сентября 2026 года}

\head{Theoretical minimum.} $GL_n(\mathbb Z)$ and unimodular matrices; the Smith normal form and
the index of a sublattice; $|GL_n(\mathbb F_q)|$ and Gaussian binomial coefficients; reduction
modulo a prime; the Cauchy--Binet formula and the matrix-tree theorem; Frobenius' theorem on the
eigenvalues of an integer matrix of finite order.

\par\vspace{-1pt}\noindent\rule{\linewidth}{0.4pt}\par\vspace{7pt}

\begin{enumerate}
\item Compute $|GL_n(\mathbb F_q)|$ and the number of $k$-dimensional subspaces of
$\mathbb F_q^n$. Let $X$ be the number of fixed vectors of a uniformly random element of
$GL_n(\mathbb F_q)$. Prove that for $n\geqslant k$ the moment $\mathsf EX^k$ does not depend on
$n$ and equals the total number of subspaces of $\mathbb F_q^k$; in particular, find
$\mathsf EX$ and the variance of $X$ for $n\geqslant2$.

\item (\foreignlanguage{russian}{ФКН, осень} 2021) Let $A$ be a $2\times2$ matrix with rational
entries and suppose that $A^n=-E$ for some positive integer $n$. Prove that $A^2=-E$ or
$A^3=-E$.

\item (SEEMOUS 2009) Let $SL_2(\mathbb Z)$ be the set of integer $2\times2$ matrices of
determinant $1$.
\begin{enumerate}[label=\alph*),align=left]
\item Find $A,B,C\in SL_2(\mathbb Z)$ with $A^2+B^2=C^2$.
\item Show that there are no $A,B,C\in SL_2(\mathbb Z)$ with $A^4+B^4=C^4$.
\end{enumerate}

\item (VJIMC 2018, cat.~I) Find all primes $p$ such that $p^3$ divides the determinant of the
$(p+6)\times(p+6)$ matrix whose diagonal entries are $2^2,3^2,4^2,\dots,(p+7)^2$ and all of whose
off-diagonal entries are equal to $1$.

\item (Romania national olympiad 1997) Let $m\geqslant2$, $n\geqslant1$ and let $A=(a_{ij})$ be an
integer $n\times n$ matrix. Prove that for every permutation $\sigma\in S_n$ there is a function
$\varepsilon\colon\{1,\dots,n\}\to\{0,1\}$ such that replacing the entries
$a_{\sigma(1)1},\dots,a_{\sigma(n)n}$ by $a_{\sigma(i)i}+\varepsilon(i)$ yields a matrix whose
determinant is not divisible by $m$.

\item (IMC) For a prime $p$ let $GL_2(\mathbb Z/p\mathbb Z)$ be the group of invertible $2\times2$
matrices over $\mathbb Z/p\mathbb Z$ and let $S_p$ be the symmetric group on $p$ elements. Show
that there is no injective group homomorphism $\varphi\colon GL_2(\mathbb Z/p\mathbb Z)\to S_p$.

\item (IMC) Ivan writes the matrix $\begin{pmatrix}2&3\\2&4\end{pmatrix}$ on the board and repeats
the following operation: he chooses a row or a column and multiplies or divides it entry-wise by
the other row or column respectively. Can he end up with
$\begin{pmatrix}2&4\\2&3\end{pmatrix}$?

\item (\foreignlanguage{russian}{Мехмат}, 2015) Find all primes $p$ for which the following is
true: if a $2015\times2015$ matrix over the field with $p$ elements has the property that its
rank drops whenever an arbitrary nonzero entry is replaced by zero, then every row contains at
most one nonzero entry.

\item (SEEMOUS 2012) Let $A=(a_{ij})$ be the $n\times n$ matrix in which $a_{ij}$ is the
remainder of $i^j+j^i$ upon division by $3$. Find the greatest $n$ for which $\det A\ne0$.

\item (Kirchhoff) Prove the matrix-tree theorem: the number of spanning trees of a graph equals
any cofactor of its Laplacian. Then compute the Smith normal form of the Laplacian of $K_n$ and
interpret the answer as the structure of the critical group.
\end{enumerate}
\end{document}
